\chapter{Introduction} \label{chap:intro}

The first chapter provides an overview of the critical role software logging plays in modern
software development and maintenance, highlighting how effective logging supports system 
reliability, debugging, and monitoring. It also addresses the significant challenges posed 
by inconsistent logging practices and heterogeneous log data, which can hinder both manual 
and automated analysis. By establishing the practical and research relevance of these issues, 
the chapter sets the stage for the subsequent examination of logging guidelines and the need 
for systematic approaches to improve log quality.
In response to these challenges, this chapter first outlines the motivation for addressing 
inconsistencies in software logging, and then articulates the specific research problem that arises from these issues.
Then, the next chapter describes the goals and research questions on logging practices and guidelines.
The structure of this thesis is outlined in the following sections.

\section{Motivation}
Modern software systems generate large amounts of log data, providing valuable insights into runtime behaviour.
Logs are widely used for debugging, system monitoring, anomaly detection, failure diagnosis and root cause analysis.
In large-scale, distributed software systems, the volume of generated log data can become too large for efficient manual inspection.
As a result, automated log analysis techniques have become very important for supporting developers and
operators in analysing system behaviour and identifying relevant events \cite{dai2022logram,zhu2019toolsandbenchmarksforautomatedlogparsing}.
The effectiveness of automated log analysis does not only depend on the analysis technique.
It also depends on the quality and consistency of the log data that is provided as input.
The latest studies show that such inconsistencies can reduce accuracy and reliability \cite{khan2022guideliensforassessingtheaccuracy,gholamian2022comprehensivesurveyloggingsoftware,chen2021asurveyofsoftwareloginstrumentation}.
Automated techniques generally require log messages to contain useful information and to follow structures to process reliably.
Inconsistent message formats, missing contextual information, unsuitable log levels or poorly placed log statements can make the analysis harder
and may reduce the quality of the generated log data.
\\
This creates an important problem because individual developers already influence log quality.
Developers have to decide what information should be logged, where logging statements should be placed,
how log messages should be created and which log level should be used \cite{he2021asurveyonautomatedloganalysis,chen2017characterizinganddetectingantipatterns}.
Even in industry, logs remain inconsistent \cite{he2022anempiricalsutdyofloganalysisatmicrosoft}.
Developers often make these decisions during software development, and they may differ across developers, teams, projects, programming languages, and application domains.
This makes investigating logging guidelines an interesting research problem.
It remains unclear which logging recommendations are provided, how specific 
and detailed these recommendations are, and what extent different sources offer similar guidance, 
and which aspects of logging receive the greatest attention.
Moreover, it is also relevant to investigate whether guidelines exist for the different decisions developers face when implementing logging
and whether these recommendations could support generating logs better suited to automated analysis. 
\\
Consequently, even software systems that generate large amounts of log data may not generate logs well suited for later automated analysis.
Existing research has investigated different approaches for improving automated log analysis,
for example, through more effective parsing and analysis techniques.
However, improving the analysis method alone does not solve problems that already exist in the input data.
If important information is missing or logging statements are implemented inconsistently,
an analysis technique can only process the information originally recorded.
Logging practices and recommendations are documented in open-source projects \cite{yuan2012characterzingloggingpractices} and in grey literature \cite{url_medium_logging_best_practices_2022,url_loggly_logging_scale_2022}.
However, this guidance is distributed across different sources, 
formats, and software contexts, making it difficult to obtain a structured overview of the available recommendations. 
Therefore, it is also relevant to investigate how logging practices and recommendations used during software development
can improve log data at its source.
\\
Observations from industrial software projects show that
logging structures and conventions can be inconsistent or only partially documented,
supporting this thesis's motivation.
Especially in larger projects involving several developers and software components, incomplete
or unclear logging guidance can result in different logging styles within the same system.
A more structured understanding of existing logging guidelines could therefore support more consistent logging practices
and provide a foundation for improving the quality of generated log data.
\\
Motivated by this, this thesis investigates existing logging guidelines and practices
and structures them into a catalogue based on relevant logging decisions.

\section{Problem Statement}
Automated log analysis techniques are widely used to help developers understand and maintain
software systems. These techniques are commonly applied for tasks such as debugging, system
monitoring, anomaly detection and root cause analysis. Most automated log analysis methods 
assume that log data are structured, consistent, and contain enough information to support analysis.
In practice, this assumption often does not hold.
\\
Developers usually write log messages manually, often without clear or standardised
logging guidelines. As a result, log formats and message contents vary across software projects,
development teams and programming languages. While some log messages contain detailed and
structured information, others lack important context or follow inconsistent formats. This inconsistency
complicates automated tools and reduces the reliability of log analysis.
Although many existing studies focus on improving log analysis techniques, such as log parsing
or anomaly detection, less attention is given to how log data is created in the first place. The
quality of the input logs is often taken for granted. When log data are inconsistent, incomplete or
poorly structured, even advanced automated analysis techniques may fail to produce accurate and
meaningful results.
\\
Logging guidelines and practices can improve log quality by specifying
what to log, when to log, and how to structure log messages. In
practice, however, developers often rely on personal experience or make logging decisions on the
spot without a systematic and reusable approach. Logging rules are rarely documented in a clearly and
consistently and are often spread across documentation and source code. This makes it difficult to
identify common practices or apply them consistently across different projects and programming
languages.
\\
As a result, developers can not refer to a single structured reference that brings
together the different logging recommendations documented across academic and
practical sources. Without a systematic approach to collecting and organising
logging guidance, it is difficult to identify recurring recommendations, differences
between sources, and areas where guidance is limited.
\\
A structured catalogue of logging guidelines can therefore provide a consolidated
reference for investigating existing logging recommendations and support more
consistent consideration of logging decisions during software development.


\section{Goals and Research Questions}

The goal of this thesis is to develop a structured and traceable catalogue of
logging guidelines and to investigate the logging decisions represented within
practical software-development guidance. The catalogue covers four central logging
dimensions: what to log, where to log, how to log, and which log levels or verbosity
to use.
\\
To achieve this goal, the study combines three complementary perspectives. First,
the study investigates developer practices and challenges through an industry survey.
Second, the study analyses academic literature to identify existing logging recommendations
and establish an initial literature-derived classification. Third, 
the study collects logging guidance from open-source repositories and transforms it into a
consolidated guideline catalogue.
\\
The academic recommendation set is then compared with the consolidated
open-source catalogue to examine how research recommendations are
represented in the collected practical guidance. The catalogue's thematic structure
is analysed using a deductive--inductive classification in which
literature-derived themes provide the initial structure, and additional themes are
introduced where required by the open-source guideline content.
\\
Therefore, this thesis formulates the research questions, presented in Table \ref{table:t1rq}.
\begin{table}[H]
\centering
\begin{tabular}{|c | p{0.6\linewidth}|}
 \hline
 \textbf{RQ1} & What logging practices and challenges do software developers encounter in software development and maintenance? \\
 \hline
 \textbf{RQ2} & What logging guidelines are documented in academic literature? \\
 \hline
 \textbf{RQ3} & What logging guidelines are documented in open-source repositories? \\
 \hline
 \textbf{RQ4} & What themes can be identified within the logging guideline catalogue? \\
 \hline
\end{tabular}
\caption{Research Questions}
\label{table:t1rq}
\end{table}


\hyperref[table:t1rq]{\textbf{RQ1}} investigates the logging practices and challenges encountered by software developers during software development and maintenance.
Insights into logging practices and challenges show how software developers struggle to implement consistent logging statements.
\\
This study investigates industry practices through a survey of software developers. The survey focuses on log implementation, log analysis, and 
fault localisation practices.
Understanding these practices and challenges from real-world industry provides a reference for existing and missing logging guidelines.
\\

\hyperref[table:t1rq]{\textbf{RQ2}} focuses on identifying logging guidelines from academic literature.
Examining research and academic literature shows how far researchers have advanced in analysing logging.
\\
Most research has investigated different aspects of software logging, including log placement, log content, log levels, log quality and log analysis.
The literature includes both logging recommendations and studies that discuss empirical, conceptual, or tool-oriented evidence about logging.
This gives the catalogue a research-based perspective and shows which recommendations have already been discussed and validated in academic work.
\\

\hyperref[table:t1rq]{\textbf{RQ3}} investigates logging guidelines in open-source software repositories.
Guidelines from open-source repositories provide a practical source for logging guidance. This can be extracted from
developers' documentation, configuration files, Markdown files and other repository artefacts. 
Characteristics such as the number of commits, contributors, programming language and application field provide contextual information about these repositories.
\\
Guidelines are often distributed across repositories and may not be explicitly described as logging guidelines.
Therefore, different search queries and collection strategies are used to identify relevant results while maintaining
sufficient coverage and relevance.
High development activity and community participation can indicate that a project is actively maintained and used.
At the same time, language and field information help determine the software contexts in which particular logging guidelines occur.
This resulting collection forms the empirical basis for developing a structured catalogue of logging guidelines from open-source software repositories.
\\

\hyperref[table:t1rq]{\textbf{RQ4}} examines the themes and characteristics that result from the identified logging guidelines.
Classifying the collected guidelines by theme, pattern, and characteristics supports creating a structured, reusable guideline catalogue.
This includes aspects such as what information should be logged, where logging should occur,
how log messages should be structured and which log levels or verbosity should be used.
\\
The open-source guidelines are classified within these dimensions using the literature-derived codebook as the deductive starting point. 
When the initial themes do not sufficiently capture the collected guidance, additional themes are developed inductively from the guideline content.
This classification forms the thematic structure of the reusable logging guideline catalogue and supports analysis of the concerns represented within the collected practical guidance.


\section{Organisation of this Thesis}
Following the introduction, Chapter 2, \textit{Background}, explains logging concepts to support the subsequent chapters.
This chapter includes related work to provide context for current research.
\\
Chapter 3, \textit{Methodology}, explains how the research is investigated.
Therefore, this chapter explains in detail how logging guidelines are collected for the literature search
and open-source project.
This includes analysing academic literature, collecting and processing open-source repository guidance, 
constructing and validating the logging guideline catalogue, and classifying it thematically and comparing its coverage.
\\
Chapter 4, \textit{Results}, contains the results derived from the methodology.
It provides a detailed analysis of the collected data and the developed logging guidelines catalogue.  
\\
Finally, the \textit{Discussion and Conclusion} present the interpretation and potential use of the logging guideline catalogue.
It also addresses limitations and possible challenges that may arise in future work.